\section{Retours primitifs et fréquence relevée}
\subsection{Automate local d’extension}
Soit \(\mathscr X_\Gamma\) l’ensemble fini des états locaux compatibles avec la route. Un état conserve

\[
 x=(t,g,B,p,\varepsilon,c,w),
\]
où \(t\) est le temps discret, \(g\) la phase nonadique, \(B\) le bloc, \(p\) le préfixe, \(\varepsilon\) le feuillet, \(c\) la retenue et \(w\) l’enroulement. Il existe une arête \(x\to y\) lorsque \(y\) est une extension admissible de \(x\), respecte le registre généalogique et passe l’élagage \(G_9\). Le tableau de 36.864 compositions de retenue fournit la loi locale d’addition ; les tableaux de retour fixent la mémoire.

Soit \(U_\Gamma\) l’opérateur de transfert normalisé de ce graphe et \(R_\Gamma\) l’involution d’orientation. On exige

\[
 U_{C_M\Gamma}=J U_\Gamma J^{-1},\qquad
 R_{C_M\Gamma}=-J R_\Gamma J^{-1}.
\]
\Needspace{7\baselineskip}
\subsection{Dénombrements primitifs orientés}
Pour \(n\ge1\), la trace orientée

\[
 a_n(\Gamma)=\operatorname{Tr}(R_\Gamma U_\Gamma^n)
\]
dénombre les retours fermés avec signe. L’inversion de Möbius sépare les orbites primitives :

\[
 p_n(\Gamma)
 =\frac1n\sum_{d\mid n}\mu(d)a_{n/d}(\Gamma).
\]
On prend l’électron comme origine affine des corrections relatives et on définit

\[
 \boxed{
 \eta_h(\Gamma)=p_h(\Gamma)-p_h(\Gamma_e),
 \qquad h\in\mathfrak S.}
\]
On obtient ainsi l’application

\[
 \mathcal F_{\rm tower}:
 \Gamma^{\rm enr}\longmapsto
 \eta(\Gamma)=(\eta_h(\Gamma))_{h\in\mathfrak S}.
\]
Il n’y a pas de paramètres par espèce. Le même automate, la même involution et la même inversion de Möbius agissent sur les 324 routes.

\Needspace{7\baselineskip}
\subsection{Convergence du raffinement primitif}
Soit \(N_\Gamma=\dim\mathscr X_\Gamma\). Si \(U_\Gamma\) est une contraction,

\[
 |a_n(\Gamma)|\le N_\Gamma.
\]
Le nombre de diviseurs de \(n\) croît de manière sous-linéaire, de sorte que

\[
 |p_n(\Gamma)|
 \le\frac{N_\Gamma\tau(n)}n.
\]
Comme \(0<q_+<q_-<1\),

\[
 |s_h|\le q_-^h+q_+^h\le2q_-^h.
\]
Par comparaison,

\[
 \sum_{h\in\mathfrak S}|\eta_h(\Gamma)s_h|<\infty
\]
uniformément sur toute collection finie de fibres. Le caractère raffiné est donc bien défini.

\Needspace{7\baselineskip}
\subsection{Identités de contrôle}
La construction doit satisfaire simultanément :

\begin{enumerate}
\item invariance par renumérotation des états ;
\item covariance sous conjugaison de feuillet ;
\item compatibilité avec la somme directe de composantes ;
\item égalité lors du recalcul à partir du registre généalogique et du graphe certifié ;
\item conservation de l’origine \(\eta(\Gamma_e)=0\) ;
\item insensibilité absolue à une permutation des colonnes externes de masse ;
\item détermination des queues au moyen d’une borne antérieure à la confrontation.
\end{enumerate}
La représentation exacte d’un résiduel donné démontre la complétude du codomaine, mais ne remplace pas ces sept preuves.

\Needspace{7\baselineskip}
\section{Horloge de route et normalisation absolue}
\Needspace{7\baselineskip}
\subsection{Les retours entiers ne suffisent pas à distinguer les espèces}
Dans l’atlas de 324 routes, le premier retour de cellule est 27, le retour cinématique au sein de CT--108 est 54 et le retour cinématique étendu est 216. Étant communs, ces entiers fixent l’architecture temporelle globale, mais ne peuvent pas expliquer à eux seuls les différences de masse entre espèces.

\Needspace{7\baselineskip}
\subsection{Holonomie temporelle}
L’information spécifique réside dans l’holonomie de l’opérateur d’extension. Soit \(\lambda_j(\Gamma)\) une valeur propre non triviale de \(U_\Gamma\) sur le sous-espace persistant. Nous écrivons

\[
 \lambda_j(\Gamma)=r_j(\Gamma)e^{i\theta_j(\Gamma)}.
\]
La fréquence interne associée est

\[
 \omega_j(\Gamma)
 =\frac{\theta_j(\Gamma)+2\pi w_j(\Gamma)}{108t_0},
\]
où l’entier \(w_j\) est fixé par l’enroulement enregistré, non par un choix de branche ultérieur. Pour un mode stable \(r_j=1\) ; pour une résonance \(r_j<1\) et son atténuation apporte une largeur.

Le quotient

\[
 \rho_{\omega,j}(\Gamma)
 =\frac{\omega_j(\Gamma)}{\omega_0}
\]
est adimensionnel et peut modifier les rapports de masse sans altérer la décade commune \(-34\). L’énergie et la masse du mode sont

\[
 E_j(\Gamma)=\hbar_{\rm ret}\omega_j(\Gamma)
 \chi_{\rm full}(\Gamma)R_{\rm act}^{\kappa_j(\Gamma)},
\]
\[
 m_j(\Gamma)=E_j(\Gamma)c_{\rm int}^{-2}.
\]
Cette équation réunit action, horloge, caractère et lecteur sans transformer aucun d’eux en un autre. Son application numérique exige de déterminer \(U_\Gamma\), de fixer l’enroulement et de déclarer la coordinatisation dimensionnelle.
